\chapter{GRUB Bootloader}
\label{ch:grub}

GRUB (Grand Unified Bootloader) version 2.x in BIOS mode (i386-pc) is
responsible for loading the Veilbox kernel from the VEILBOX ext4 partition.
This chapter explains the GRUB boot process, the blocklist format, and
the rootless installation procedure.

\section{GRUB Architecture}

GRUB 2.x for BIOS systems consists of three distinct stages:

\subsection{Stage 1: boot.img}

\texttt{boot.img} is a 512-byte file written to LBA 0 (the Master Boot
Record of the disk). It contains:

\begin{itemize}[nosep]
  \item A 446-byte real-mode assembly stub that sets up segment registers,
        initializes the stack pointer, and reads the blocklist from a fixed
        offset within the core.img file.
  \item The 64-byte MBR partition table at offset 0x1BE–0x1FD.
  \item The 2-byte boot signature (0x55 0xAA) at offset 0x1FE–0x1FF.
\end{itemize}

When the BIOS (SeaBIOS) jumps to address 0x7C00, \texttt{boot.img}
executes. It reads the blocklist pointer embedded in \texttt{core.img}
to determine how many sectors to load and where.

\subsection{Stage 1.5: core.img}

\texttt{core.img} is written to the sectors immediately following the MBR
(starting at LBA 1). Its structure is:

\begin{figure}[H]
\centering
\begin{tikzpicture}[node distance=0.5cm]
\node[box, text width=6cm] (diskboot) {diskboot.img (LBA 1, 1 sector)};
\node[box, text width=6cm, below of=diskboot] (kernel2) {GRUB kernel \& filesystem drivers};
\node[box, text width=6cm, below of=kernel2] (cfg) {Embedded config string};
\draw[arrow] (diskboot) -- node[right] {Loads} (kernel2);
\draw[arrow] (kernel2) -- node[right] {Executes} (cfg);
\end{tikzpicture}
\caption{Structure of core.img on disk. diskboot.img loads the GRUB
kernel and its filesystem drivers, which then parse the embedded config.}
\label{fig:coreimg}
\end{figure}

The first sector (\texttt{diskboot.img}) is loaded by \texttt{boot.img}.
It switches the CPU from real mode to protected mode, enables the A20
gate (allowing access to memory above 1 MB), and loads the rest of
\texttt{core.img}.

The GRUB kernel embedded in core.img provides:
\begin{itemize}[nosep]
  \item Command-line parsing (the GRUB shell)
  - Device management (disk detection, partition parsing)
  - Filesystem drivers (ext2, part\_msdos, biosdisk)
  - Module loading infrastructure
\end{itemize}

The embedded configuration (added by \texttt{grub2-mkimage -c}) tells
GRUB to find the VEILBOX partition and set the root prefix:

\begin{lstlisting}
search --label --set=root VEILBOX
set prefix=($root)/boot/grub
\end{lstlisting}

\subsection{Stage 2: /boot/grub/}

Once core.img finds the VEILBOX partition and loads the ext2 filesystem
driver, it reads \texttt{/boot/grub/grub.cfg}. The default Veilbox
configuration:

\begin{lstlisting}
set default=0
set timeout=5

menuentry "Veilbox" {
    linux /boot/vmlinuz console=ttyS0 console=tty1
}
\end{lstlisting}

\section{The Blocklist Format}

core.img contains a blocklist at offset 0x1F4 that tells diskboot.img
how many sectors to load. The format is:

\begin{table}[H]
\centering
\begin{tabular}{|c|c|c|c|}
\hline
\textbf{Offset} & \textbf{Size} & \textbf{Field} & \textbf{Description} \\
\hline
0x1F4 & 4 bytes & start & First LBA sector number \\
0x1F8 & 4 bytes & start\_hi & Must be 0 (high bits) \\
0x1FC & 2 bytes & remaining & Number of sectors to load \\
0x1FE & 2 bytes & seg & Real-mode segment address \\
\hline
\end{tabular}
\caption{core.img blocklist format at offset 0x1F4. Total: 12 bytes.}
\label{tab:blocklist}
\end{table}

\section{Rootless GRUB Installation}

Installing GRUB traditionally requires loop-mounting the disk image and
running \texttt{grub2-install}, both of which need root privileges.
Veilbox uses a Python script that writes directly to the raw disk image
file.

\subsection{Installation Procedure}

\begin{enumerate}[nosep]
  \item Extract the partition table from the disk image using sfdisk.
  \item Write the 512-byte \texttt{boot.img} to LBA 0.
  \item Restore the saved partition table to bytes 0x1BE–0x1FD.
  \item Generate a fresh \texttt{core.img} using \texttt{grub2-mkimage}
        with the embedded config.
  \item Write \texttt{core.img} starting at LBA 1.
  \item Copy GRUB modules (\texttt{*.mod}) and configuration file to
        \texttt{/boot/grub/} on the VEILBOX partition using debugfs.
\end{enumerate}

\subsection{Why Not Patch the Blocklist}

Earlier versions of the installer attempted to patch the blocklist in
core.img to reflect the actual size of the core image. This approach
failed because:

\begin{itemize}[nosep]
  \item The blocklist is 12 bytes, not 16 as previously assumed.
        Patching byte 8 (intended as the \texttt{seg} field) actually
        overwrites \texttt{start\_hi} with the sector count, corrupting
        the blocklist.
  \item \texttt{grub2-mkimage} already generates a correct blocklist
        with \texttt{start=2}, \texttt{start\_hi=0},
        \texttt{remaining=N}, and \texttt{seg=0x0820}.
  \item Any modification invalidates the GRUB geometry check at
        \texttt{diskboot.img:0x63}, causing a ``geometry error'' triple
        fault.
\end{itemize}

\subsection{The Python Installer}

\begin{lstlisting}[language=python]
def install_grub(disk, core_img, boot_img):
    with open(disk, 'r+b') as f:
        # Save partition table (bytes 0x1BE-0x1FD)
        f.seek(0x1BE)
        pt = f.read(64)

        # Write boot.img to LBA 0
        f.seek(0)
        f.write(boot_img)

        # Restore partition table
        f.seek(0x1BE)
        f.write(pt)

        # Write core.img starting at LBA 1
        f.seek(512)
        f.write(core_img)
\end{lstlisting}

\section{GRUB Interactive Shell}

GRUB provides an interactive shell for troubleshooting and manual booting.
Access it by pressing `c` at the GRUB menu:

\begin{lstlisting}
grub> ls
(hd0) (hd0,msdos1)

grub> ls (hd0,msdos1)/
boot/

grub> ls (hd0,msdos1)/boot/
vmlinuz grub/

grub> set root=(hd0,msdos1)
grub> linux /boot/vmlinuz console=ttyS0
grub> boot
\end{lstlisting}

\subsection{Useful GRUB Commands}

\begin{longtable}{|p{3cm}|p{10cm}|}
\hline
\textbf{Command} & \textbf{Purpose} \\
\hline
\texttt{ls} & List devices and partitions. \\
\texttt{ls \textless device\textgreater/} & List files on a partition. \\
\texttt{set} & Display all variables. \\
\texttt{set root=...} & Set root partition. \\
\texttt{linux} & Load a Linux kernel. \\
\texttt{initrd} & Load an initrd (not needed with embedded initramfs). \\
\texttt{boot} & Boot the loaded kernel. \\
\texttt{reboot} & Reboot the system. \\
\texttt{configfile} & Load a configuration file. \\
\texttt{search} & Search for a partition by label or UUID. \\
\texttt{insmod} & Load a GRUB module. \\
\texttt{cat} & Display file contents. \\
\hline
\end{longtable}

\section{GRUB Configuration Advanced Topics}

\subsection{Multiple Boot Entries}

The grub.cfg can define multiple boot options:

\begin{lstlisting}
set default=0
set timeout=5

menuentry "Veilbox (Standard)" {
    linux /boot/vmlinuz console=ttyS0 console=tty1
}

menuentry "Veilbox (Debug)" {
    linux /boot/vmlinuz console=ttyS0 console=tty1 \
        debug ignore_loglevel earlyprintk=serial
}

menuentry "Veilbox (Single User)" {
    linux /boot/vmlinuz console=ttyS0 console=tty1 single
}

menuentry "Reboot" {
    reboot
}
\end{lstlisting}

\subsection{Serial Console in GRUB}

GRUB can output to the serial console, which is useful for headless systems:

\begin{lstlisting}
serial --unit=0 --speed=38400 --word=8 --parity=no --stop=1
terminal_input serial console
terminal_output serial console

menuentry "Veilbox" {
    linux /boot/vmlinuz console=ttyS0 console=tty1
}
\end{lstlisting}

\section{GRUB Module Reference}

The modules listed during \texttt{grub2-mkimage} control what features
GRUB has access to:

\begin{longtable}{|p{4cm}|p{8cm}|}
\hline
\textbf{Module} & \textbf{Purpose} \\
\hline
ext2 & Read ext2/3/4 filesystems (needed for VEILBOX partition). \\
part\_msdos & Parse MBR partition tables. \\
biosdisk & Access disks via BIOS INT 13h. \\
search & Search for partitions by label, UUID, or file. \\
configfile & Read GRUB configuration files. \\
serial & Serial console I/O. \\
terminal & Terminal management (serial + VGA). \\
ls & List devices and directories. \\
cat & Display file contents. \\
linux & Load Linux kernel. \\
boot & Execute loaded kernel. \\
reboot & Reboot the system. \\
halt & Halt the system. \\
\hline
\end{longtable}
